\section*{第 \ref{ch:g2:air-around-us} 章 --- 身边的空气}

\begin{solution}{exo:g2:air-around-us:1}
例如：风吹弯树木或吹得树叶沙沙作响；
风筝或旗帜在空中飘扬；寒冷的早晨呼出的白气；
用吸管吹气时冒出的泡泡；穿堂风把门吹得砰的一声关上。
\end{solution}

\begin{solution}{exo:g2:air-around-us:2}
里面困住了空气。挤压时想把空气压进更小的空间，
空气便会反推回来——真正空的东西则不会。
\end{solution}

\begin{solution}{exo:g2:air-around-us:3}
杯子里装满了空气，被困住的空气占据着自己的空间：
两样东西不能同时占据同一个地方，
所以水无法升到纸巾那里。
\end{solution}

\begin{solution}{exo:g2:air-around-us:4}
不要倾斜杯子。倾斜后，被困住的空气会以大气泡的形式向上逃走，
水便占据了它的位置——纸巾也就湿透了。
\end{solution}

\begin{solution}{exo:g2:air-around-us:5}
每一声“咕嘟”，都是一份空气以气泡的形式逸出，
同时有一口水补进它的位置。两者轮流进出，是因为瓶颈狭窄，
空气与水无法同时通过，
它们也不能同时占据瓶内的同一个空间。
\end{solution}

\begin{solution}{exo:g2:air-around-us:6}
空气——一直装到瓶口。“空”只是指“饮料已喝完”；
瓶子里其实满是看不见的东西。
\end{solution}

\begin{solution}{exo:g2:air-around-us:7}
风就是流动的空气。空气的作用有：鼓起船帆、
放风筝、晾干衣服、转动风车玩具——任选两种。
\end{solution}

\begin{solution}{exo:g2:air-around-us:8}
空气从气球口向后喷出，
气球便向相反的方向向前飞去。
\end{solution}

\begin{solution}{exo:g2:air-around-us:9}
瓶子里已经装满了空气，而密封的漏斗使空气无路可逃——所以水无法进入。
解决办法：松开密封（或把漏斗稍微抬起），
让空气从旁边逸出，
水便能欢快地“咕嘟咕嘟”流下去。
\end{solution}

\begin{solution}{exo:g2:air-around-us:10}
是的！把装满空气的杯子倾斜，使其杯口位于另一只杯子杯口之下：空气以一个个银色的大气泡逸出，
\emph{向上}进入装满水的杯子，并聚集在顶部，
把水向下挤出。空气上升、水下落，
是因为空气总会穿过水向上浮起——
这正是颠倒过来的倒水。
\end{solution}
